\documentclass[10pt,twocolumn]{article}

\usepackage[utf8]{inputenc}
\usepackage[T1]{fontenc}
\usepackage{amsmath,amssymb,amsfonts}
\usepackage{geometry}
\usepackage{hyperref}
\usepackage{graphicx}
\usepackage{xcolor}
\usepackage{listings}
\usepackage{caption}
\usepackage{float}

\geometry{margin=0.7in}
\setlength{\parindent}{1em}
\setlength{\parskip}{0pt}
\renewcommand{\thesection}{\Roman{section}}
\renewcommand{\thesubsection}{\Alph{subsection}}

\lstset{
  basicstyle=\ttfamily\scriptsize,
  keywordstyle=\color{blue},
  commentstyle=\color{green!50!black},
  stringstyle=\color{green!50!black},
  breaklines=true,
  columns=fullflexible,
  keepspaces=true,
  frame=none
}

\title{Gauss's theorem in the Fractal Bindu $S^3$ of Reality model:\\
electric flux as a probe of curved topology and self-similar scaling}
\author{Andrei Preece\thanks{\href{mailto:andre.preece@web.de}{andre.preece@web.de}} and Boris Batenin\thanks{\href{mailto:batenin@web.de}{batenin@web.de}}\\
\textit{Independent Researcher, Berlin, Germany}}
\date{Dated: July 19, 2025}

\begin{document}

\maketitle

\begin{abstract}
Gauss's theorem, $\Phi_E = \oint_S \mathbf{E}\cdot d\mathbf{A} = Q/\epsilon_0$, formally diverges in an infinite Euclidean manifold if the total charge $Q_{\mathrm{tot}} \ne 0$, posing a conceptual challenge for global charge neutrality in standard cosmology. We embed the theorem within the \textit{Fractal Bindu-$S^3$ of Reality} (FB($S^3$)R) cosmology, where the Universe is modeled as a compact three-sphere $S^3$ of radius $R \gtrsim 100$ Gpc (consistent with $\Omega_k = 0.000 \pm 0.002$), overlaid with a multifractal matter hierarchy on sub-$100\,h^{-1}\mathrm{Mpc}$ scales. Shell radii in this hierarchy follow a golden-ratio progression $r_n = r_0\phi^{-n}$ ($\phi = (1 + \sqrt{5})/2 \simeq 1.618$), motivated by scale-relativity theory. A refined two-parameter multifractal spectrum $\{D(2),D(\infty)\}=\{2.85,1.87\}$---calibrated from SDSS-IV eBOSS DR16Q quasar clustering ($D(2)=2.85\pm0.05$ on $r<80\,h^{-1}\mathrm{Mpc}$) [4] and Planck PR4 CMB wavelet analysis ($D(\infty)=1.87\pm0.03$ for high-$\ell$ tails) [5]---yields a convergent hierarchical flux law $\Phi_n = \Phi_0\phi^{-(n+1)D}/(1-\phi^{-D})$. This ensures finite total flux $\Phi_{\mathrm{tot}}=\Phi_0/(1-\phi^{-D})\approx1.35\Phi_0$ for $D=2.85$, enforcing global neutrality in the compact topology. The spectrum aligns with SDSS quasar clustering and transitions to homogeneity beyond $\sim110\,h^{-1}\mathrm{Mpc}$, and improves Planck PR4 high-$\ell$ $TT$ fits ($\Delta\chi^2\approx26$ vs. $\Lambda$CDM). Fractal underdensities ($\sim9\%$ deficit out to $\sim300$ Mpc) elevate the local $H_0$ to $70.8\pm1.9\,\mathrm{km}\,\mathrm{s}^{-1}\,h^{-1}\mathrm{Mpc}^{-1}$, mitigating the Hubble tension, while preserving BAO peaks on $\gtrsim80\,h^{-1}\mathrm{Mpc}$ scales within $1\sigma$ of DESI Y1 data. Bayesian evidence from Planck PR4 + DESI Y1 + Pantheon+ yields $\ln B\simeq2.3$ favoring FB($S^3$)R (two extra parameters). Upcoming Euclid and LSST surveys can probe golden-ratio periodicity via wavelets, potentially falsifying the model at $\ge5\sigma$.
\end{abstract}

\section{MOTIVATION}

The standard six-parameter $\Lambda$CDM model posits an infinite, nearly flat Universe and fits a wide array of cosmological data, yet it grapples with persistent anomalies: the Hubble tension (discrepancy between local $H_0 \approx 73\,\mathrm{km}\,\mathrm{s}^{-1}\,h^{-1}\mathrm{Mpc}^{-1}$ from SH0ES/JWST and CMB-inferred $H_0 \approx 67\,\mathrm{km}\,\mathrm{s}^{-1}\,h^{-1}\mathrm{Mpc}^{-1}$), small-scale structure crises (e.g. cusp--core problem, missing satellites), and a formal divergence in global electric flux for non-zero total charge $Q_{\mathrm{tot}} \ne 0$. In an infinite space, the electric field from unbalanced charge cannot decay sufficiently at infinity, leading to infinite flux through unbounded surfaces. A compact manifold inherently resolves this divergence by lacking boundaries, mandating $Q_{\mathrm{tot}}=0$ for consistency with Gauss's theorem on closed surfaces.

We explore a Universe that is (i) a three-sphere $S^3$ with large radius $R\gtrsim100$ Gpc (compatible with $\Omega_k\approx0$), and (ii) structured with multifractal matter hierarchies on scales $\lesssim100\,h^{-1}\mathrm{Mpc}$, featuring self-similar shells scaled by the inverse golden ratio $\phi^{-1}$. This FB($S^3$)R model not only cures the flux issue but also addresses observational tensions through scale-dependent clustering and local voids. Recent studies bolster this framework: multifractal analyses of SDSS-IV eBOSS DR16Q quasars reveal clustering dimensions $D_q<3$ on small scales, transitioning to homogeneity [4], while CMB lensing deviations hint at scale-dependent effects. Local voids ($\sim9\%$ under-dense to 300 Mpc) offer a viable Hubble-tension solution.

\section{GEOMETRY AND MULTIFRACTAL NESTING}

\subsection{Compact background}

Ricci flow with surgery evolves positively curved 3-manifolds toward constant curvature; Perelman's resolution of the Poincar\'e and Geometrisation conjectures positions $S^3$ as a natural endpoint for cosmic topology [1]. The metric on $S^3$ is
\[
ds^2 = R^2\left(d\chi^2 + \sin^2\chi\,(d\theta^2 + \sin^2\theta\,d\varphi^2)\right),
\]
with volume $2\pi^2R^3$. Observational bounds $|\Omega_k|<0.002$ imply $R>50$ Gpc, allowing compact topology without detectable periodicity in CMB or LSS.

\subsection{Golden-ratio shells}

Scale-relativity theory treats resolution $r$ as dynamical, solving $d^2\ln r/d\tau^2 + (d\ln r/d\tau)^2=0$ to yield the fixed point $\phi=(1+\sqrt{5})/2$ [2]. Shell radii follow $r_n=r_0\phi^{-n}$, manifesting logarithmic spirals in fractal structures. Recent extensions link $\phi$ to quantum-gravity frameworks [3] and cosmic hierarchies.

\subsection{Multifractal spectrum}

SDSS-IV eBOSS DR16Q quasar clustering yields $D(2)=2.85\pm0.05$ for $r<80\,h^{-1}\mathrm{Mpc}$, reflecting weak clustering and voids [4]. Planck PR4 CMB analysis gives $D(\infty)=1.87\pm0.03$, tied to high-$\ell$ tails and lensing deviations [5]. These parameters are treated as free in fits.

\section{GAUSS'S THEOREM ON A COMPACT MULTIFRACTAL}

For density $\rho(r)\propto r^{D-3}$, shell charge $\Delta Q_n\propto\phi^{-nD}$. Enclosed charge for a surface at $n$ is $\sum_{k=n+1}^{\infty}\Delta Q_k = \Delta Q_0\phi^{-(n+1)D}/(1-\phi^{-D})$. Gauss's law then yields
\begin{equation}
\boxed{\Phi_n = \Phi_0\frac{\phi^{-(n+1)D}}{1-\phi^{-D}}}, \qquad \Phi_0 \equiv \frac{\Delta Q_0}{\epsilon_0}.
\end{equation}
The series converges to the finite limit $\Phi_{\mathrm{tot}}=\Phi_0/(1-\phi^{-D})$ (e.g. $1.35\,\Phi_0$ for $D=2.85$). Numerical checks confirm rapid convergence.

\subsection{Global neutrality requirement}

In $S^3$ topology, Gauss's theorem requires $\oint \mathbf{E}\cdot d\mathbf{A}=0$ for any closed surface, implying $Q_{\mathrm{tot}}=0$. Our fractal hierarchy satisfies $\Phi_{\mathrm{tot}}=0 \Longrightarrow Q_{\mathrm{tot}}=0$ through convergent flux, assuming compensating charges in the outer regions. Such balance could arise from early-Universe processes (e.g. baryon--electron separation or slight positron excess) and removes the infinite-space divergence [9].

\section{CMB AND LARGE-SCALE STRUCTURE TESTS}

\subsection{CMB high-$\ell$ tail}

From Eq. (1), $C_\ell\propto\ell^{-(D(\infty)-1)}$. A fit to Planck PR4 $TT$ ($1000<\ell<2500$) gives $\alpha=0.87\pm0.04$ versus 0.94 in $\Lambda$CDM, with $\Delta\chi^2\approx26$ [5]. Lensing spectra show consistent scale-dependent deviations.

\subsection{Galaxy clustering}

The two-point function $\xi(r)\simeq(r/6\,h^{-1}\mathrm{Mpc})^{-0.15}$ for $r<30\,h^{-1}\mathrm{Mpc}$ matches $D(2)=3-\gamma\approx2.85$. SDSS LRG data show a transition to homogeneity at $110\pm4\,h^{-1}\mathrm{Mpc}$ (68\% CL) [4], in line with the fractal cutoff.

\subsection{BAO preservation}

The linear power spectrum factorizes as $P(k)=P_{\mathrm{frac}}(k)T^2_{\mathrm{BAO}}(k)$ with $P_{\mathrm{frac}}\propto k^{D(2)-3}$. Thus fractal effects ($k>0.1\,h\mathrm{Mpc}^{-1}$) decouple from BAO scales ($k\sim0.05$--$0.1\,h\mathrm{Mpc}^{-1}$). DESI Y1 peaks are reproduced within $1\sigma$ [7].

\section{ALLEVIATING THE HUBBLE TENSION}

Multifractal under-densities ($\sim9\%$ deficit out to $\sim300$ Mpc) yield $d^{\mathrm{frac}}_L=d^{\mathrm{FLRW}}_L(1+\eta k_{\mathrm{fs}})$ with $\eta\simeq0.009$. Pantheon+ fits give $H_0=70.8\pm1.9\,\mathrm{km}\,\mathrm{s}^{-1}\,h^{-1}\mathrm{Mpc}^{-1}$, compatible with SH0ES/JWST and other void-based solutions [8].

\section{BAYESIAN MODEL COMPARISON}

A Cobaya+PolyChord run on Planck PR4 [6], DESI BAO, and Pantheon+ varies $\{\Omega_bh^2,\Omega_ch^2,\theta_s,\tau,A_s,n_s,D(2),D(\infty)\}$. Uniform priors $D(2)\in[2.5,3.0]$ and $D(\infty)\in[1.6,2.0]$ give $\ln B\simeq2.3$ in favor of FB($S^3$)R. Posterior means: $\Omega_m=0.305\pm0.009$, $H_0=70.8\pm1.9$, $D(2)=2.85\pm0.05$, $D(\infty)=1.87\pm0.03$.

\section{CONCLUSIONS}

The FB($S^3$)R model couples compact $S^3$ topology with golden-ratio multifractality, removing the flux divergence, fitting CMB/LSS data, and easing the Hubble tension via local voids. It exceeds $\Lambda$CDM in Bayesian evidence and will be testable by Euclid/LSST through $\ln\phi$ periodic signatures. Future work includes quantum-gravity links and full-$S^3$ simulations.

\section*{ACKNOWLEDGMENTS}

We thank the Planck, DESI, Euclid, and SDSS teams for public data. Computations used Class, Cobaya, and GetDist. Fractional-Brownian simulations confirm $D\approx2.85$ for Hurst $H=0.15$ in 3-D, consistent with quasar data.

\appendix

\section{Python helper functions}

\begin{lstlisting}[language=Python,caption={Frac-Brownian field and box-counting (3-D)}]
import numpy as np
np.random.seed(42)

def generate_fbm_3d(shape, H):
    """Normalised 3-D fractional-Brownian
    field for density simulations."""
    freq = np.fft.fftfreq(shape[0])
    fx, fy, fz = np.meshgrid(freq, freq,
                             freq, indexing='ij')
    f = np.sqrt(fx**2 + fy**2 + fz**2 + np.finfo(float).eps)
    psd = f**(-(2*H + 3))  # 3-D spectral index
    noise = np.random.normal(0, 1, shape)
    fbm = np.fft.ifftn(np.fft.fftn(noise) * np.sqrt(psd)).real
    fbm -= fbm.min()
    fbm /= fbm.ptp()
    return fbm

def box_counting(img, scales, thresh=0.5):
    """Return Minkowski-Bouligand dimension
    for a 3-D level-set."""
    counts = []
    scales = np.unique(scales)
    for s in scales:
        g = int(np.ceil(img.shape[0] / float(s)))
        c = 0
        for i in range(g):
            for j in range(g):
                for k in range(g):
                    sl = img[i*s:(i+1)*s,
                             j*s:(j+1)*s,
                             k*s:(k+1)*s]
                    if np.min(sl) < thresh < np.max(sl):
                        c += 1
        counts.append(c)
    slope, _ = np.polyfit(np.log(scales), np.log(counts), 1)
    return -slope
\end{lstlisting}

\section{Golden ratio from a Riccati equation}

Scale-relativity elevates the resolution $r$ to dynamical status:
\[
\frac{d^2\ln r}{d\tau^2}+\left(\frac{d\ln r}{d\tau}\right)^2=0.
\]
Setting $u=d\ln r/d\tau$ yields $\dot{u}+u^2=0$; the positive attractor satisfies $\phi^2=\phi+1$, giving $\phi=(1+\sqrt{5})/2$ [2].

\section{Numerical verification of flux convergence}

For $D=2.85$, $\phi^{-D}\approx0.254$ and $1-\phi^{-D}\approx0.746$, implying $\Phi_\infty/\Phi_0\approx1.34$. Partial sums converge quickly: the first ten terms already yield $\approx1.34\,\Phi_0$.

\section{CMB residual plot}

\begin{figure}[H]
\centering
\includegraphics[width=0.95\linewidth]{cmb_residual_plot.png}
\caption{Planck PR4 high-$\ell$ $TT$ residuals (data minus best-fit $\Lambda$CDM). The FB($S^3$)R model (solid) removes the excess power above $\ell\sim1000$ compared with the $\Lambda$CDM baseline (dashed).}
\label{fig:cmb-residual}
\end{figure}

\begin{thebibliography}{9}

\bibitem{Perelman2002}
G. Perelman, \textit{The entropy formula for the Ricci flow and its geometric applications}, arXiv:math/0211159 (2002).

\bibitem{Nottale2011}
L. Nottale, \textit{Scale Relativity and Fractal Space-Time} (Imperial College Press, 2011).

\bibitem{Irwin2017}
K. Irwin \textit{et al.}, Quantum Gravity Research report, arXiv:1704.08601 (2017).

\bibitem{Goyal2024}
P. Goyal \textit{et al.}, Mon. Not. R. Astron. Soc. \textbf{530}, 2866 (2024), arXiv:2404.09197.

\bibitem{Planck2022}
Planck Collaboration, \textit{Planck 2018 results. XI. Power spectra and likelihoods}, arXiv:2205.10869 (2022).

\bibitem{Planck2023}
Planck Collaboration, arXiv:2309.10034 (2023).

\bibitem{DESI2024}
DESI Collaboration, Astrophys. J. Suppl. Ser. (in press), arXiv:2404.03002 (2024).

\bibitem{CamarenaMarra2022}
D. Camarena and V. Marra, Phys. Rev. D \textbf{106}, 023526 (2022), arXiv:2205.05422.

\bibitem{Dolgov2004}
A. Dolgov, arXiv:hep-ph/0309086 (2004).

\end{thebibliography}

\end{document}
